\chapter{Conclusiones y trabajos futuros}\label{chap:conclusiones}

\section{Conclusiones}

Este trabajo muestra que el valor de la señal metabolómica no reside solo en discriminar AD y NC con un rendimiento razonable, sino en cómo obliga a replantear el problema al llevar esa discriminación a una decisión clínica real. La aportación principal no es, por tanto, ``un clasificador binario mejor'', sino una formulación de triaje que hace explícita la incertidumbre cuando la evidencia deja de ser suficiente.

Los resultados sostienen con claridad que la cohorte contiene señal biológica real, pero también que esa señal no admite una lectura simplista. Los tamaños de efecto univariantes, la estructura multivariante del problema y el comportamiento consistente del modelo base indican que la metabolómica plasmática residualizada aporta información útil para diferenciar AD y NC. Al mismo tiempo, el solapamiento entre clases, el bajo EPV y la distribución de la varianza en PCA muestran que la señal está distribuida, no concentrada en un biomarcador dominante, y que los modelos muy flexibles pagan rápido un coste en varianza. En este contexto, la combinación de \textit{stability selection}, \textit{forward selection}, ingeniería de variables y regresión logística regularizada permitió construir un panel de 11 variables que mejora al conjunto original tanto en rendimiento como en control del optimismo.

Esa parte, sin embargo, no agota el problema clínico. El modelo base logra un rendimiento suficiente para justificar el panel metabolómico, pero no resuelve de forma homogénea a todos los pacientes. Los errores más persistentes no se comportan como ruido simple ni como outliers eliminables: apuntan a una ambigüedad biológica real desde la perspectiva metabolómica periférica. Por eso, seguir optimizando una clasificación plana con umbral único habría sido metodológicamente pobre. La cuestión no era solo mejorar una métrica, sino representar con honestidad dónde el modelo dispone de evidencia fuerte y dónde deja de tenerla.

Desde ahí se justifica la segunda aportación del trabajo: la reformulación del problema como un sistema de dos etapas con activación condicional del segundo escalón. El Stage 1 resume la evidencia bioquímica y define zonas de decisión directa y de ambigüedad; el Stage 2 reinterpreta esa ambigüedad con un bloque clínico ortogonal. Bajo esta formulación, la evaluación multi-seed muestra una mejora clara frente al baseline metabolómico directo, con \textbf{BA = 0,910 $\pm$ 0,051}, \textbf{sensibilidad = 0,943 $\pm$ 0,056} y un \textbf{gap train-test cercano a 1,2 puntos porcentuales}. Además, las alternativas más complejas exploradas ---blending, stacking, correcciones por vecindario o ampliaciones del Stage 2--- no sostuvieron mejoras robustas, reforzando que la solución final surge de comparación empírica, no de arbitrariedad.

Esa conclusión exige una lectura crítica. El sistema retenido no debe presentarse como una solución universal ni como la versión conceptualmente más pura de una arquitectura secuencial. El contraste con la logística monolítica deja claro que, si el único criterio fuese maximizar BA global, en esta cohorte existen alternativas mejores. Del mismo modo, la variante más estricta (entrenar el segundo escalón solo sobre la zona amarilla) tampoco mejoró el rendimiento medio: en un experimento de 100 particiones externas con búsqueda interna de $C$ sobre la zona amarilla, esa configuración obtuvo BA = 0,898 $\pm$ 0,044 frente a BA = 0,910 $\pm$ 0,051 de la configuración final, siendo superior en apenas 13 de las 100 semillas. La elección final debe leerse como un compromiso pragmático: no es la formulación más pura ni la que optimiza una métrica aislada, pero sí la que ofrece el mejor equilibrio observado entre sensibilidad alta, explicitación de incertidumbre y estabilidad empírica dentro de las limitaciones del dataset.

Es importante subrayar que el sistema de dos etapas no maximiza BA global. Una regresión logística monolítica que combine directamente metabolómica y variables clínicas alcanza BA = 0,930 ± 0,045, comparado con BA = 0,910 ± 0,051 del sistema propuesto. El trabajo conscientemente intercambia esos 2 puntos porcentuales de BA por tres ganancias clínicas: sensibilidad más alta (0,943 vs 0,904) que reduce falsos negativos; una zona explícita de ambigüedad que el clínico puede ver y gestionar; y una estructura interpretable donde metabolómica y clínica aparecen por separado, facilitando la auditoría del razonamiento del sistema. Esta decisión refleja una postura metodológica deliberada: la utilidad clínica no se reduce a una única métrica numérica, y la transparencia sobre la incertidumbre tiene un valor igual o mayor que la optimización métrica pura.

En suma, la aportación central del trabajo no es maximizar BA, sino demostrar que la incertidumbre debe modelarse explícitamente. El trabajo no está solo en construir un clasificador metabolómico aceptable, sino en demostrar que ese clasificador, por sí solo, no describe bien el problema clínico. Lo relevante es la reformulación: pasar de una clasificación plana a un sistema de dos etapas más coherente con un escenario de priorización, donde la utilidad depende tanto de acertar como de reconocer cuándo la evidencia todavía no basta.

Este TFM no propone simplemente un clasificador metabolómico, sino un modo más honesto de representar la incertidumbre diagnóstica en la enfermedad de Alzheimer mediante un sistema de priorización clínica escalonada.

\section{Limitaciones}

Los resultados deben interpretarse dentro de límites claros. El primero es muestral: 148 sujetos siguen siendo pocos para afirmar transportabilidad amplia o justificar arquitecturas más complejas sin reservas. El segundo es de validación: las conclusiones proceden de validación interna rigurosa, pero descansan en una única cohorte, por lo que la generalización a otros entornos clínicos o preanalíticos no puede darse por garantizada. A ello se añade un límite de diseño: el trabajo distingue AD establecida frente a NC, pero no aborda aún el problema longitudinal de detectar conversión futura o estadios prodrómicos.

Existen además limitaciones específicas del pipeline. Una parte relevante del valor del panel descansa en la residualización previa, por lo que cambios en ese proceso pueden modificar la señal disponible. La cobertura directa del Stage 1, cuando se exige seguridad estricta en las zonas extremas, sigue siendo necesariamente modesta. Y permanece una tensión conceptual que conviene explicitar: la arquitectura final se describe mejor como un sistema de dos etapas con activación condicional que como una cascada estricta entrenada solo sobre la subpoblación amarilla. Esta precisión no debilita el trabajo; evita atribuirle una pureza metodológica que la implementación final no pretende.

\section{Líneas de trabajo futuro}

Las líneas de continuidad más relevantes no pasan por añadir complejidad algorítmica de forma ciega, sino por ampliar la evidencia de la formulación propuesta. La prioridad inmediata es validar externamente el panel y el sistema de dos etapas en cohortes independientes, y extenderlo después a escenarios clínicos más exigentes (deterioro cognitivo leve, sujetos preclínicos o circuitos asistenciales reales), donde la utilidad se mida no solo en BA, sino también en tiempos diagnósticos, pruebas evitadas y seguridad observada.

En un plano más biológico, el siguiente paso natural sería combinar la firma metabolómica con otros biomarcadores proteicos y genéticos ---más allá del APOE ya incorporado en el Stage 2--- para evaluar si la misma lógica de dos etapas sigue siendo la estructura más útil en un espacio de información más rico. Y, quizá más importante, conviene profundizar en la caracterización del subgrupo ambiguo. Los pacientes frontera son los que más tensionan el sistema y, al mismo tiempo, los que pueden aportar más información sobre los límites biológicos reales de la señal disponible.